\documentclass[10pt]{article}
\usepackage[T1]{fontenc}
\usepackage{lmodern,microtype,amsmath,amssymb,amsthm,booktabs}
\usepackage[margin=0.85in]{geometry}
\usepackage[colorlinks=true,linkcolor=blue,urlcolor=blue,citecolor=blue]{hyperref}
\usepackage{bookmark}
\newtheorem{theorem}{Theorem}
\newtheorem{lemma}[theorem]{Lemma}
\newtheorem{corollary}[theorem]{Corollary}
\newcommand{\dd}{\mathrm d}
\newcommand{\sech}{\operatorname{sech}}
\newcommand{\artanh}{\operatorname{artanh}}
\setlength{\emergencystretch}{2em}
\title{Positive Moment Bounds and Growing Plunge Windows}
\author{Kirtus Justus (Theoretical Physicist)\\\small Port Coquitlam, British Columbia, Canada}
\date{October 2026}
\hypersetup{pdftitle={Positive Moment Bounds and Growing Plunge Windows},pdfauthor={Kirtus Justus (Theoretical Physicist)}}
\begin{document}
\maketitle
\begin{abstract}
The reciprocal trace associated with the even half-filled sine kernel has a
positive Stieltjes representation. It follows that every finite inverse-power
expansion has a remainder of the sign of, and smaller than, its first omitted
term. Two rational bounds control the exact central slope. The full formal
inverse-power series nevertheless has zero radius of convergence. An elementary
variational construction controls growing central windows of the commuting
Jacobi matrix. An exact Darboux deletion, with its singular endpoint domains
proved explicitly, realizes the omitted spectrum as a positive differential
operator. Elementary bracketing then supplies higher reciprocal moments and
proves the signed plunge law uniformly for $|\lambda|\le L_N=o(N^{1/4})$.
In particular a logarithmic number of actual central pairs obey the law with
relative error $O(N^{-1/2})$. No bulk quantum-period remainder is assumed.
\end{abstract}

\section{The finite problem and exact inputs}
Let $N=2m$ be a positive even integer. With site indices $1,\ldots,N$, define
\[
B_{jk}=\begin{cases}1/2,&j=k,\\
\sin[\pi(j-k)/2]/[\pi(j-k)],&j\ne k,
\end{cases}\qquad
(J_N)_{j,j+1}=(J_N)_{j+1,j}=b_j:=\frac{j(N-j)}N,
\quad (J_N)_{jj}=0.
\]
The matrices commute and $0<B<I$. Write their common eigenvalues as
$J_Nv=\lambda v$ and $Bv=\zeta v$. Let $H=NJ_N/2$,
$\alpha=(N^2-1)/4$, and let $\mu_{N,k}$, in increasing order starting at
$k=0$, be the eigenvalues of the Friedrichs realization
\[
A_Nu=-(\cos x\,u')'-\alpha\cos x\,u,
\qquad -\pi/2<x<\pi/2.
\]
The endpoint condition excludes the logarithmic solution. The exact
spectral-tail and trace identities established in \cite{tails,trace} are
\begin{gather}
\{\mu_{N,0},\ldots,\mu_{N,N-1}\}=-\operatorname{spec}H,
\qquad \mu_{N,N}>\|H\|,\qquad
Z_N:=\sum_{k\ge N}\mu_{N,k}^{-1},\label{eq:tail}\\
\log\frac\zeta{1-\zeta}=F_N(N\lambda/2),\qquad
F_N(b)=2\sum_{k\ge N}\artanh\frac b{\mu_{N,k}},\label{eq:logit}\\
Z_N=\pi\int_0^\infty e^{-Nt}f(t)\dd t,
\qquad f(t)=P_{-1/2}(\cosh t)^2.\label{eq:laplace}
\end{gather}
Here $P$ is the real Legendre function on $[1,\infty)$. The first identity
identifies precisely the modes removed from the differential spectrum; the
second holds at the actual common finite modes. The normalization in
\eqref{eq:laplace} follows from a weighted Legendre resolvent and a Mellin
reflection identity, rather than an asymptotic fit. These previous results
are the exact inputs to this paper. The positive integral in
\eqref{eq:laplace} defines an interpolation for every real $N>0$; equality
with the specified spectral trace is used only for even integers.

\section{A positive Fourier measure}
\begin{lemma}[Autocorrelation identity]
Put $g(u)=\sech(u)^{1/2}$ and use the Fourier convention
$\widehat g(\omega)=\int_{\mathbb R}e^{-i\omega u}g(u)\dd u$. Then
\begin{align}
y(t):=P_{-1/2}(\cosh t)
&=\frac1\pi\int_{\mathbb R}g(u)g(u-t)\dd u,\label{eq:correlation}\\
\widehat g(\omega)&=\frac{|\Gamma(1/4+i\omega/2)|^2}{\sqrt{2\pi}},
\qquad
y(t)=\int_{\mathbb R}e^{i\omega t}p(\omega)\dd\omega,\label{eq:density}\\
p(\omega)&=\frac{|\Gamma(1/4+i\omega/2)|^4}{4\pi^3}.
\end{align}
The function $p$ is an even, strictly positive probability density with
moments of every order. Consequently $f=y^2$ is the Fourier transform of
the probability density $q=p*p$.
\end{lemma}
\begin{proof}
The classical angular representation \cite[Eq.~14.12.7]{dlmf} gives
\[
y(t)=\frac1\pi\int_0^\pi
(\cosh t+\sinh t\cos\theta)^{-1/2}\dd\theta.
\]
For $t\ge0$, substitute $\tan(\theta/2)=e^u$. Then
$\dd\theta=\sech u\,\dd u$ and
\[
\cosh t+\sinh t\cos\theta=\frac{\cosh(u-t)}{\cosh u},
\]
which proves \eqref{eq:correlation}. Both sides are even in $t$, so it
holds on the whole real line. Substitution $x=e^{2u}$ in the Fourier
integral, followed by Euler's beta integral \cite[Eq.~5.12.3]{dlmf}, gives
\[
\widehat g(\omega)=2^{-1/2}
\int_0^\infty x^{1/4-i\omega/2-1}(1+x)^{-1/2}\dd x
=\frac{\Gamma(1/4-i\omega/2)\Gamma(1/4+i\omega/2)}{\sqrt{2\pi}}.
\]
The function $g$ and all its derivatives decay exponentially; in
particular $g$ is a Schwartz function. Fourier inversion for its
autocorrelation proves \eqref{eq:density}. Plancherel and
$\int_{\mathbb R}g^2\dd u=\pi$ give $\int p=1$. Positivity follows
because the gamma function has no zeros. The Schwartz property supplies
every moment and justifies differentiation under the Fourier integral.
Multiplication of the two absolutely convergent Fourier integrals gives
the convolution statement.
\end{proof}

Let $\sigma$ be the image of $q(\omega)\dd\omega$ under $s=\omega^2$.
It is a probability measure on $[0,\infty)$, has full support, and has
moments $m_k=\int_0^\infty s^k\dd\sigma(s)$ for all $k\ge0$.
\begin{theorem}[Stieltjes representation and signed remainder]
For the interpolation in \eqref{eq:laplace}, for every $N>0$,
\begin{equation}
\boxed{\quad\frac{NZ_N}\pi=\int_0^\infty
\frac{\dd\sigma(s)}{1+s/N^2}.\quad}\label{eq:Stieltjes}
\end{equation}
For every integer $K\ge0$,
\begin{equation}
0<(-1)^{K+1}\left[\frac{NZ_N}\pi-
\sum_{k=0}^{K}\frac{(-1)^km_k}{N^{2k}}\right]
<\frac{m_{K+1}}{N^{2K+2}}.\label{eq:nextterm}
\end{equation}
The formal series in $N^{-2}$ has zero radius of convergence, while
\eqref{eq:nextterm} proves its asymptotic validity to every fixed order.
\end{theorem}
\begin{proof}
Insert the Fourier representation for $f$ in \eqref{eq:laplace}.
Absolute Fubini is justified by
$\int_{\mathbb R}q(\omega)\int_0^\infty Ne^{-Nt}\dd t\,\dd\omega=1$.
Integrating in $t$ and using the evenness of $q$ gives
$N^2/(N^2+\omega^2)$, proving \eqref{eq:Stieltjes}. The finite identity
\[
\frac1{1+x}=\sum_{k=0}^{K}(-x)^k+
\frac{(-1)^{K+1}x^{K+1}}{1+x},\qquad x\ge0,
\]
then gives the signed remainder in \eqref{eq:nextterm} exactly as
\[
\frac1{N^{2K+2}}\int_0^\infty
\frac{s^{K+1}}{1+s/N^2}\dd\sigma(s).
\]
It is strictly between zero and $m_{K+1}/N^{2K+2}$ because $\sigma$ is
not supported at zero. For every $R>0$,
$m_k\ge R^k\sigma([R,\infty))$, with a strictly positive last factor.
Therefore $m_k^{1/k}\to\infty$, proving zero convergence radius.
The estimate at each fixed $K$ already proves the asymptotic assertion.
\end{proof}

Since $f^{(2k)}(0)=(-1)^km_k$, its differential equation or
hypergeometric series determines the moments exactly:
\begin{equation}
\begin{array}{c|rrrrrr}
k&0&1&2&3&4&5\\\hline
m_k&1&\frac14&\frac{17}{64}&\frac{169}{256}&
\frac{50777}{16384}&\frac{1575169}{65536}.
\end{array}\label{eq:moments}
\end{equation}
For example, if
\[
S_N=\frac\pi N-\frac\pi{4N^3}+\frac{17\pi}{64N^5}
-\frac{169\pi}{256N^7}+\frac{50777\pi}{16384N^9},
\]
then for every positive even $N$,
\begin{equation}
0<S_N-Z_N<\frac{1575169\pi}{65536N^{11}}.\label{eq:concrete}
\end{equation}
These are bounds on the exact trace, not just on formal coefficients.

\begin{corollary}[Two rational slope bounds]
For every positive even $N$,
\begin{equation}
\boxed{\quad\frac{4\pi}{16N^2+17}<\pi-NZ_N<
\frac\pi{4N^2+1}.\quad}\label{eq:rational}
\end{equation}
Also
\begin{equation}
\left|\pi-NZ_N-\frac\pi{4N^2}\right|
<\frac{17\pi}{64N^4}.\label{eq:slope}
\end{equation}
\end{corollary}
\begin{proof}
Strict Jensen applied to $s\mapsto(1+s/N^2)^{-1}$ under $\sigma$
gives $NZ_N/\pi>(1+m_1/N^2)^{-1}$. For the probability measure
$\dd\tau=s\dd\sigma/m_1$, its mean is $m_2/m_1=17/16$ and
\[
1-\frac{NZ_N}\pi=\frac{m_1}{N^2}
\int\frac{\dd\tau(s)}{1+s/N^2}
>\frac{m_1}{N^2+m_2/m_1}.
\]
Both measures are nondegenerate, so both inequalities are strict.
Substitution of the moments proves \eqref{eq:rational}; the case
$K=1$ of \eqref{eq:nextterm} proves \eqref{eq:slope}.
\end{proof}

\section{A growing window of small Jacobi eigenvalues}
List the positive eigenvalues of $J_{2m}$ increasingly as
$0<\lambda_{1,+}\le\cdots\le\lambda_{m,+}$. The negative eigenvalues
are their opposites: conjugation by alternating signs sends $J$ to $-J$.
There is no zero eigenvalue, since
$\det J_{2m}=(-1)^m\prod_{k=1}^m b_{2k-1}^2\ne0$.
\begin{theorem}[Explicit logarithmic-window bound]
Set $L=\lfloor m/2\rfloor$. For any integer $r\ge1$ with $16^r\le L$,
put $\ell=\lfloor\log L/(r\log16)\rfloor$. Then
\begin{equation}
\lambda_{r,+}\le\frac{12\sqrt2\,\pi}{\ell\log16}.
\label{eq:packetbound}
\end{equation}
In particular, if $1\le r\le\log L/(2\log16)$, then
\begin{equation}
\boxed{\quad\lambda_{r,+}\le
\frac{24\sqrt2\,\pi r}{\log L}.\quad}\label{eq:logbound}
\end{equation}
Thus the innermost $r_N$ pairs shrink uniformly to zero whenever
$r_N=o(\log N)$ as even $N\to\infty$.
\end{theorem}
\begin{proof}
Define $A_1=1$ and $A_{k+1}=-b_{2k-1}A_k/b_{2k}$ for $k<m$.
For $v_{2k-1}=A_kf_k$, $v_{2k}=0$, direct substitution gives
\begin{equation}
\|v\|^2=\sum_{k=1}^m A_k^2f_k^2,\qquad
\|Jv\|^2=\sum_{k=1}^m(b_{2k-1}A_k)^2(f_k-f_{k+1})^2,
\label{eq:gauge}
\end{equation}
where $f_{m+1}=0$ encodes the actual last row. For $1\le k\le L$,
\begin{equation}
\frac1{4k}\le A_k^2\le\frac3k,
\qquad (b_{2k-1}A_k)^2\le12k.\label{eq:weights}
\end{equation}
Indeed, with $c_s=(1/2)_s/s!$,
\[
|A_k|=c_{k-1}\prod_{j=1}^{k-1}\left(1+\frac1{2(m-j)}\right).
\]
For $s\ge1$, $s c_s^2$ increases from $1/4$, because its successive
ratio is $(2s+1)^2/[4s(s+1)]>1$. Also $(s+1)c_s^2$ decreases from one,
because its successive ratio is $1-(3s+2)/[4(s+1)^3]<1$.
These prove the required lower bound and $c_{k-1}^2\le1/k$.
For $k\le L$, the logarithm of the squared product is at most
\[
\sum_{j=1}^{k-1}\frac1{m-j}\le\frac{k-1}{m-k+1}<1.
\]
The squared product is therefore below $e<3$. Checking $k=1$ directly and using
$b_{2k-1}\le2k$ proves \eqref{eq:weights}.

Set $R=16^\ell$ and $h=\log R$. On each integer block
$[a,b]=[R^j,R^{j+1}]$, $j=0,\ldots,r-1$, use the sequence
\[
f_k=\begin{cases}
\sin[\pi\log(k/a)/h],&a\le k\le b,\\
0,&\text{otherwise}.
\end{cases}
\]
All blocks lie in $[1,L]$, and their shared endpoints are zero.
The resulting odd-site vectors have disjoint nonzero supports; so do
their images under $J$, since a block contributes only to the even
rows corresponding to $k=a,\ldots,b-1$ in \eqref{eq:gauge}.
The norm and energy forms consequently have no cross terms between
different packets.

The numbers $aR^{1/4}$ and $aR^{3/4}$ are integers, and $f_k^2\ge1/2$
between them. Thus
\[
\|v\|^2\ge\frac18\sum_{k=aR^{1/4}}^{aR^{3/4}}\frac1k
\ge\frac18\int_{aR^{1/4}}^{aR^{3/4}+1}\frac{\dd x}{x}
\ge\frac h{16}.
\]
The sine function is 1-Lipschitz and $\log(1+1/k)\le1/k$, so
$|f_{k+1}-f_k|\le\pi/(hk)$ inside the block. Therefore
\[
\|Jv\|^2\le\frac{12\pi^2}{h^2}\sum_{k=a}^{b-1}\frac1k
\le\frac{12\pi^2(h+1)}{h^2}\le\frac{18\pi^2}{h},
\]
where $h\ge\log16>2$. Every nonzero vector in the $r$-dimensional
packet span thus has $\|Jv\|^2/\|v\|^2\le288\pi^2/h^2$.

Order the odd sites before the even sites. Then
\[
J=\begin{pmatrix}0&T^{\mathsf T}\\T&0\end{pmatrix},
\]
where $T$ is invertible and upper bidiagonal, with diagonal
$b_1,b_3,\ldots,b_{2m-1}$. Its singular values are precisely the
positive eigenvalues of $J$. Min--max on the odd-site block
$T^{\mathsf T}T$ therefore gives
$\lambda_{r,+}^2\le288\pi^2/h^2$, proving \eqref{eq:packetbound}.
If $\log L/(r\log16)\ge2$, the floor is at least half its argument;
this proves \eqref{eq:logbound}. Finally $\log L\sim\log N$.
\end{proof}
The constants are deliberately elementary, not optimal. The finite-size
condition for \eqref{eq:packetbound} already requires $N\ge64$ when
$r=1$; the assertion is not applied outside its stated range.

\section{Uniform transfer to the modular error}
For a common mode define
\[
\eta(t)=\pi t\,{}_3F_2(1/4,1/2,3/4;1,3/2;4t^2),\qquad
\varepsilon=\log\frac{1-\zeta}{\zeta},\qquad
e=\varepsilon+N\eta(\lambda/N).
\]
All quantities are dimensionless. In particular the word ``modular''
here does not supply a physical clock or an energy unit.
\begin{lemma}[Explicit ratio control]
For every even $N\ge8$ and every common mode with $0<|\lambda|\le1$,
\begin{equation}
\left|\frac{-4N^2e}{\varepsilon}-1\right|
\le\frac{10}{N^2}+33\lambda^2.\label{eq:ratio}
\end{equation}
\end{lemma}
\begin{proof}
The two central sites give $\|H\|\ge N^2/8$ by a Rayleigh quotient.
From \eqref{eq:tail}, with $\beta=N\lambda/2$,
\[
\frac{|\beta|}{\mu_{N,N}}\le\frac{4|\lambda|}N\le\frac12,
\qquad
\sum_{k\ge N}\mu_{N,k}^{-3}\le\frac{64Z_N}{N^4}
\le\frac{64\pi}{N^5}.
\]
The positive artanh series therefore yields
\begin{equation}
|F_N(\beta)-N\lambda Z_N|
\le\frac{64\pi|\lambda|^3}{9N^2}.\label{eq:cubic}
\end{equation}
The coefficients $d_k$ of the hypergeometric series in the definition
of $\eta$ are positive and decreasing, with $d_1=1/16$. Indeed their
successive ratio is
$[(k+1/4)(k+1/2)(k+3/4)]/[(k+1)^2(k+3/2)]<1$.
Thus
\[
|N\eta(\lambda/N)-\pi\lambda|
\le\frac{\pi|\lambda|^3}{4N^2(1-4\lambda^2/N^2)}
\le\frac{4\pi|\lambda|^3}{15N^2}.
\]
Since $\varepsilon=-F_N(\beta)$, combining this with
\eqref{eq:cubic} gives
\[
|e-\lambda(\pi-NZ_N)|\le\frac{8\pi|\lambda|^3}{N^2}.
\]
Set $a=4N^2e/(\pi\lambda)$ and
$b=F_N(\beta)/(\pi\lambda)$. Equations \eqref{eq:slope} and
\eqref{eq:cubic}, including the common sign of each artanh term and
$\lambda$, imply
\[
|a-1|\le\frac{17}{16N^2}+32\lambda^2,\qquad
b\ge1-\frac1{4N^2}>0,\qquad |b-1|\le\frac8{N^2}.
\]
Consequently
\[
\left|\frac ab-1\right|
\le\frac{145/(16N^2)+32\lambda^2}{1-1/(4N^2)}
\le\frac{10}{N^2}+33\lambda^2.
\]
Here $a/b=-4N^2e/\varepsilon$. Strict positivity of $b$ also
justifies division by $\varepsilon$.
\end{proof}

\begin{corollary}[Uniform signed plunge law on a growing window]
Let $r_N\ge1$ be integers with $r_N=o(\log N)$ as even $N\to\infty$.
For both signs of the innermost $r_N$ pairs of common eigenmodes,
\begin{equation}
\sup_{1\le j\le r_N,\,\pm}
\left|\frac{-4N^2 e_{j,\pm}}{\varepsilon_{j,\pm}}-1\right|
=O\!\left(N^{-2}+\frac{r_N^2}{\log^2 N}\right)\longrightarrow0.
\label{eq:windowlaw}
\end{equation}
For all sufficiently large $N$, an explicit upper bound for the left
side is
\[
\frac{10}{N^2}+33\left(
\frac{24\sqrt2\,\pi r_N}{\log\lfloor N/4\rfloor}\right)^2.
\]
\end{corollary}
\begin{proof}
The bound \eqref{eq:logbound} puts every eigenvalue in the window below one
in absolute value for all sufficiently large $N$. Apply
\eqref{eq:ratio} to each mode. Its constants are uniform, so taking
the supremum preserves the estimate.
\end{proof}
In particular every fixed number of central pairs obeys the law, with
the conservative relative error $O((\log N)^{-2})$.

\section{An exact positive operator for the differential tail}

The higher reciprocal moments admit a useful positive realization.
Write $p(x)=\cos x$ and introduce the Liouville coordinate
\[
z(x)=\int_{-\pi/2}^x\frac{\dd u}{\sqrt{\cos u}},\qquad
0<z<\mathcal L,\qquad
\mathcal L=\frac{\sqrt\pi\,\Gamma(1/4)}{\Gamma(3/4)}.
\]
Below $p(z)=\cos x(z)$.
\begin{theorem}[Exact deletion of the finite block]
The spectrum of
\begin{equation}
\boxed{\quad T_N=-D_z^2+
\frac{4N^2-1}{16p(z)}+\frac{p(z)}{16}\quad}
\label{eq:positivepartner}
\end{equation}
with its Friedrichs realization is precisely
$\{\mu_{N,k}:k\ge N\}$. Its form domain is $H_0^1(0,\mathcal L)$,
and its endpoints are limit point. In particular,
\begin{equation}
\mu_{N,N}>N^2/4,\qquad
Z_N(s):=\sum_{k\ge N}\mu_{N,k}^{-s}
=\operatorname{Tr}T_N^{-s}\quad(s>1/2).
\label{eq:highertrace}
\end{equation}
\end{theorem}
\begin{proof}
The unitary map $\psi(z)=p(x(z))^{1/4}u(x(z))$ transforms $A_N$
into $-D_z^2+V_N$, where direct differentiation gives
\[
V_N=-N^2p/4+p/16-1/(16p).
\]
At either endpoint, with $d$ the distance to that endpoint,
$p=d^2/4+O(d^6)$. Thus $V_N=-1/(4d^2)+O(d^2)$, and the
selected eigenfunctions behave as $d^{1/2}$ times a regular power
series in $d^2$. The logarithmic branch is excluded. This critical
endpoint requires a domain argument in addition to the usual formal
Darboux--Crum transformation \cite{crum}.

For one deletion, let $g>0$ be the current ground eigenfunction,
with eigenvalue $E_0$, and suppose its selected endpoint exponent is
$k+1/2$, $k\ge0$. Define $Q$ as the closure from $C_c^\infty$ of
$D_z-g'/g$. The compact-support ground-state identity, followed by
closure of quadratic forms, gives
\[
S-E_0=Q^*Q
\]
for the current Friedrichs operator $S$. The self-adjoint partner
$QQ^*+E_0$ extends the differential expression with potential
$V-2(\log g)''$. Differentiable Frobenius expansions show that its
endpoint coefficient is $((k+1)^2-1/4)d^{-2}+O(1)$.
Its two branches have exponents $k+3/2$ and $-k-1/2$; only the
first is square integrable. Hence the partner is limit point at
both endpoints and is its unique self-adjoint realization.
Moreover $\ker Q^*=\{0\}$, since $Q^*f=0$ implies $f=c/g$,
which is not square integrable, including at $k=0$.
The gap above $E_0$ makes $\operatorname{ran}Q$ closed; its orthogonal
complement is $\ker Q^*$, so it is all of $L^2$.
Polar decomposition therefore pairs the orthogonal complement of
$g$ unitarily with the partner space. Exactly the ground eigenvalue
is deleted, with completeness and compact resolvent preserved.
The new ground state has exponent $k+3/2$, permitting iteration.

Direct finite Fourier transformation and \eqref{eq:tail} show that
the first $N$ original eigenfunctions span
\[
\operatorname{span}\{e^{imx}:m=-(N-1)/2,\ldots,(N-1)/2\}.
\]
Their transformed Wronskian is therefore
\[
W_z(\psi_0,\ldots,\psi_{N-1})=C\,p^{N^2/4},\qquad C\ne0.
\]
Indeed the common factor $p^{1/4}$ contributes $p^{N/4}$,
the change of variable contributes
$(\dd x/\dd z)^{N(N-1)/2}=p^{N(N-1)/4}$, and the Fourier
Wronskian is a nonzero constant Vandermonde determinant.
Iterating the one-step identity gives
$V_{\rm tail}=V_N-2D_z^2\log|W_z|$.
Since
\[
D_z^2\log p=-\tfrac12(p+p^{-1}),
\]
this is exactly \eqref{eq:positivepartner}.
The Wronskian formula evaluates the already justified self-adjoint
deletions; it does not replace their domain argument.

The final potential is $(N^2-1/4)d^{-2}+O(1)$.
Hardy's inequality gives form domain $H_0^1$.
Finally
\[
\frac{4N^2-1}{16p}+\frac p{16}-\frac{N^2}{4}
=\frac{(1-p)(4N^2-1-p)}{16p}\ge0.
\]
The kinetic energy is strictly positive on $H_0^1$, proving the
first assertion of \eqref{eq:highertrace}. The trace convergence
and its quantitative form also follow from the bracketing below.
\end{proof}

\section{Higher reciprocal moments by elementary bracketing}

\begin{theorem}[Higher-moment leading terms with an explicit error]
For each fixed $s>1/2$,
\begin{equation}
Z_N(s)=N^{1-2s}\{I_s+O_s(N^{-1/2})\},\qquad
I_s=\frac{2^{2s-2}\Gamma(s-1/2)\Gamma(s/2)}
{\Gamma(s)\Gamma((s+1)/2)}.
\label{eq:momentleading}
\end{equation}
For every integer $M\ge1$ one has, in particular,
\begin{equation}
\left|N^5Z_N(3)-3\pi\right|\le D_{N,M}:=
\frac{6\mathcal L}{M}+\frac{128M}{N}
+3\pi\left[\left(1-\frac1{4N^2}\right)^{-5/2}-1\right].
\label{eq:thirdmomentbound}
\end{equation}
\end{theorem}
\begin{proof}
Put $h=1/N$ and $Q_h=h^2T_N=-h^2D_z^2+V_h$. Then
\[
V_h=\left(1-\frac{h^2}{4}\right)V_0+\frac{h^2p}{16},
\qquad V_0=\frac1{4p},\qquad
(1-h^2/4)V_0\le V_h\le V_0.
\]
The potential $V_h$ has minimum $1/4$ and is monotone on
each half-interval. Divide each half into $M$ cells of length
$\delta=\mathcal L/(2M)$. Neumann conditions and the cell infimum
of $V_h$ give an upper trace bound by the min-max principle.
Dirichlet conditions and the cell supremum give a lower trace
bound. The two exterior cells, whose supremum is infinite, are
omitted from the latter bound; this is a restriction of variational
subspaces, not a boundedness assumption at the singular endpoints.

For constant potential $v>0$ the cell eigenvalues are
$(h\pi j/\delta)^2+v$, with $j\ge0$ for Neumann and $j\ge1$
for Dirichlet conditions. The integral test gives
\[
\left|h\operatorname{Tr}(-h^2D_{\rm N/D}^2+v)^{-s}
-c_s\delta v^{1/2-s}\right|\le hv^{-s},
\qquad c_s=\frac{\Gamma(s-1/2)}{2\sqrt\pi\,\Gamma(s)}.
\]
The total cell trace error is at most $2Mh4^s$.
The function $V_h^{1/2-s}$ rises from zero to $4^{s-1/2}$
and falls symmetrically; its upper and lower Riemann sums differ
by at most $2\delta4^{s-1/2}$. Therefore, with
$I_s(h)=c_s\int_0^{\mathcal L}V_h^{1/2-s}\dd z$,
\[
\left|h\operatorname{Tr}Q_h^{-s}-I_s(h)\right|
\le\frac{c_s\mathcal L4^{s-1/2}}{M}+2Mh4^s.
\]
This also proves convergence of the trace. Potential comparison
gives $I_s\le I_s(h)\le(1-h^2/4)^{-(s-1/2)}I_s$, where
\[
I_s=c_s4^{s-1/2}\int_{-\pi/2}^{\pi/2}
(\cos x)^{s-1}\dd x.
\]
The beta integral evaluates this as \eqref{eq:momentleading}.
Since $h\operatorname{Tr}Q_h^{-s}=N^{2s-1}Z_N(s)$,
choosing $M=\lceil\sqrt N\rceil$ proves the stated rate.
For $s=3$, $c_3=3/16$ and $I_3=3\pi$, giving
\eqref{eq:thirdmomentbound}.
\end{proof}
For example $I_5=105\pi/4$. The first-moment result obtained
by this general argument is weaker than \eqref{eq:laplace};
the new information is the independently proved higher-moment
coefficient. No bulk period expansion was used.

\section{A larger uniform eigenvalue window}

The third-moment coefficient cancels the leading nonlinear term
of the modular error.
\begin{theorem}[Uniform center law beyond shrinking eigenvalues]
Let $L_N=o(N^{1/4})$ as even $N\to\infty$. Uniformly over nonzero
common modes with $|\lambda|\le L_N$,
\begin{align}
e&=\frac{\pi\lambda}{4N^2}
\left[1+O\left(N^{-2}+\lambda^2N^{-1/2}
+\lambda^4N^{-2}\right)\right],\label{eq:largercenter}\\
\frac{-4N^2e}{\varepsilon}&\longrightarrow1.
\label{eq:largerratio}
\end{align}
This includes every fixed bounded eigenvalue window. More explicitly,
for even $N\ge2$, $0<|\lambda|\le N/4$, and any integer $M\ge1$,
\begin{equation}
\left|\frac{-4N^2e}{\varepsilon}-1\right|
\le
\frac{\displaystyle
\frac{21}{16N^2}
+\left(\frac{D_{N,M}}{3\pi}+\frac{16}{9N^2}\right)\lambda^2
+\frac{20\lambda^4}{N^2}}
{\displaystyle1-\frac1{4N^2}}.
\label{eq:largerexplicit}
\end{equation}
\end{theorem}
\begin{proof}
Set $\beta=N\lambda/2$. By \eqref{eq:highertrace},
$r=|\beta|/\mu_{N,N}<2|\lambda|/N\le1/2$.
The exact artanh series \eqref{eq:logit} gives
\[
F_N(\beta)=N\lambda Z_N+\frac{N^3\lambda^3}{12}Z_N(3)+R,
\qquad
|R|\le\frac{2|\beta|^5}{5(1-r^2)}Z_N(5).
\]
The simple bound $Z_N(5)\le(4/N^2)^4Z_N\le256\pi/N^9$
suffices, giving $|R|\le64\pi|\lambda|^5/(15N^4)$.
The convergent hypergeometric series defining $\eta$, with positive
decreasing coefficients, gives
\[
N\eta(\lambda/N)=\pi\lambda+\frac{\pi\lambda^3}{4N^2}
+R_\eta,\qquad |R_\eta|\le\frac{7\pi|\lambda|^5}{16N^4}.
\]
Equation \eqref{eq:thirdmomentbound} cancels the two cubic terms,
and \eqref{eq:slope} therefore implies
\[
\left|e-\frac{\pi\lambda}{4N^2}\right|
\le\frac{17\pi|\lambda|}{64N^4}
+\frac{D_{N,M}|\lambda|^3}{12N^2}
+\frac{5\pi|\lambda|^5}{N^4}.
\]
Choosing $M=\lceil\sqrt N\rceil$ proves \eqref{eq:largercenter}.
For division, put $a=4N^2e/(\pi\lambda)$ and
$b=F_N(\beta)/(\pi\lambda)$.
The sign of every artanh term gives
$b\ge NZ_N/\pi\ge1-1/(4N^2)>0$.
Also $Z_N(3)\le16Z_N/N^4$ and the artanh bound before the
cubic term give
\[
|b-1|\le\frac{1/4+16\lambda^2/9}{N^2}.
\]
Combining this with the preceding bound on $a-1$ gives
\eqref{eq:largerexplicit}, since $a/b=-4N^2e/\varepsilon$.
On the stated window all its terms tend uniformly to zero,
proving \eqref{eq:largerratio}.
\end{proof}
\begin{corollary}[A logarithmic number of actual central pairs]
For all sufficiently large even $N$, put
\[
r_N=\left\lfloor
\frac{\log\lfloor N/4\rfloor}{\log16}
\right\rfloor.
\]
For both signs of the innermost $r_N$ pairs of common modes,
\[
\sup_{1\le j\le r_N,\,\pm}
\left|\frac{-4N^2e_{j,\pm}}{\varepsilon_{j,\pm}}-1\right|
=O(N^{-1/2}).
\]
\end{corollary}
\begin{proof}
The choice of $r_N$ satisfies $16^{r_N}\le\lfloor N/4\rfloor$.
The integer $\ell$ in \eqref{eq:packetbound} is at least one, so
\[
|\lambda_{j,\pm}|\le\frac{12\sqrt2\,\pi}{\log16}
\qquad(1\le j\le r_N).
\]
Apply \eqref{eq:largerexplicit} on this fixed bounded window with
$M=\lceil\sqrt N\rceil$. Its constants are independent of the
selected mode. The numerical constant and the bracketing rate
are not asserted to be optimal.
\end{proof}
This is still a center result:
$t=\lambda/N=o(N^{-3/4})$. The argument proves no fixed-$t$
bulk quantum-period remainder and does not claim the displayed
bracketing rate or window exponent is optimal.

\section{Literature comparison and scope}
The logarithmic scale has a rigorous precedent. Deift, Its, and
Krasovsky \cite[Theorem 1.3]{dik} prove that, for a symbol taking two
positive constant values on complementary rational arcs, eigenvalues
in a fixed compact subinterval of the gap have spacing of order
$1/\log N$; every interval of a suitable constant times $1/\log N$
contains an eigenvalue. The choice of values $1,2$ and a half-circle
gives $I+B_N$ up to diagonal unitary conjugation. This supports the
known logarithmic occupation scale. The finite Jacobi construction
above supplies an independent bound with explicit constants. No new
scale or priority is claimed.

The positive-measure theorem controls the exact first reciprocal trace
at every even size. The growing-window result concerns actual common
eigenmodes, and its proof is uniform over the stated window. It does
not establish a uniform bulk quantum-period expansion, an odd-size
deflation identity, or a general-filling effective-size convention. The fixed-interior bulk expansion is supplied separately by \emph{Uniform Bulk Modular Asymptotics and the Leading Wilson Entropy Law}; it is not a consequence of moment bracketing alone.
In particular, the trace measure's moments $m_k$ are not the higher
reciprocal spectral traces $\sum_{j\ge N}\mu_{N,j}^{-(2k+1)}$.
Their leading asymptotics are established separately by \eqref{eq:momentleading}.

Exact arithmetic checks recover the displayed moments independently
from the hypergeometric series. Quadrature checks compare the
autocorrelation, gamma density, and Legendre representations, and
finite matrix checks verify the packet energies. Symbolic checks also verify
the transformed potential, Wronskian exponent, and moment constants. These checks provide
reproducibility; the analytic inequalities and min--max argument are
the proofs. No finite numerical agreement is used to infer a theorem.

\paragraph{Acknowledgments and provenance.}
This work belongs to \emph{From Prephysical Hypotheses to Physical Models}.
The present derivations and checks were developed by Kirtus Justus with
ChatGPT (OpenAI). The preceding connection-coefficient proposal was supplied
in correspondence prepared with Claude (Anthropic). The wider program has
also used Gemini (Google), Grok (SpaceXAI), and DeepSeek (High-Flyer); no
specific proof here is attributed to those systems. Assistance does not
imply external peer review or endorsement. No priority claim is made for
the special-function identities or the logarithmic central-mode scale.

\begin{thebibliography}{9}
\bibitem{tails} K. Justus, \emph{Spectral Tails and Exact Concentration Logits},
research manuscript, October 2026.
\bibitem{trace} K. Justus, \emph{Reciprocal Traces and the Half-Filled Plunge Law},
research manuscript, October 2026.
\bibitem{dlmf} NIST Digital Library of Mathematical Functions,
\href{https://dlmf.nist.gov/14.12.E7}{Eq.~14.12.7} and
\href{https://dlmf.nist.gov/5.12.E3}{Eq.~5.12.3}.
\bibitem{crum} M. M. Crum, Associated Sturm--Liouville systems,
\emph{Quarterly Journal of Mathematics} \textbf{6} (1955), 121--127;
\href{https://arxiv.org/abs/physics/9908019}{archival transcription}.
\bibitem{dik} P. Deift, A. Its, and I. Krasovsky,
\emph{Eigenvalues of Toeplitz matrices in the bulk of the spectrum},
\href{https://arxiv.org/abs/1110.4089}{arXiv:1110.4089}, Theorem 1.3.
\end{thebibliography}
\end{document}
